% =============================================================================
%  Cultural-life decomposition — skeleton paper
%  STATUS: UNPROVEN. Deepest test of the predictive thesis.
% =============================================================================
\documentclass[11pt]{article}

\usepackage[utf8]{inputenc}
\usepackage[T1]{fontenc}
\usepackage[expansion=false]{microtype}
\usepackage{amsmath, amssymb, amsthm, mathtools}
\usepackage[margin=1in]{geometry}
\usepackage{enumitem}
\usepackage[hidelinks]{hyperref}

\theoremstyle{plain}
\newtheorem{theorem}{Theorem}
\theoremstyle{definition}
\newtheorem{definition}[theorem]{Definition}

\newcommand{\X}{\mathbf{X}}

\title{Cultural Life at Vitality Profile $(3, (1, 1, 1))$: \\
       Decomposition of Linguistic Evolution under $\mathrm{Decomp}$ \\
       \large (UNPROVEN --- skeleton paper)}
\author{Liam McGuire\thanks{Unaffiliated researcher, Seattle, USA. liammcg44@gmail.com}}
\date{\today}

\begin{document}
\maketitle

\begin{abstract}
This paper is a \textbf{skeleton} for a worked-example paper that
specialises the decomposition functor $\mathrm{Decomp}$ on a small
slice of linguistic evolution (Great Vowel Shift, or a meme cycle on
social media) and verifies vitality profile $(3, (1, 1, 1))$ or
deeper. \textbf{No proof is attempted.} The deepest test of the
framework's predictive thesis, likely to draw productive pushback
from outside ALife. The result is UNPROVEN.
\end{abstract}

\section{Goal (UNPROVEN)}

Exhibit, for a specific linguistic-evolution slice, the four-slot
system $(\X, V, F, T) \in \mathbf{Sys}$ that the decomposition
functor produces, and verify that its vitality profile reaches at
least $(3, (1, 1, 1))$. This is the framework's strongest
substantive prediction: cultural systems are structurally deeper
than individual biological organisms.

\section{Plan of attack (no proof attempted)}

\begin{enumerate}[leftmargin=2em,itemsep=4pt]
  \item \textbf{Choose the slice (UNPROVEN).} Candidates: the Great
    Vowel Shift in English (1400--1700); a viral meme cycle (days
    to weeks); a creolisation event (decades).
  \item \textbf{Identify entities at each level (UNPROVEN).} Level 0:
    speakers. Level 1: linguistic items (words, phrases, syntactic
    structures). Level 2: languages (or major sociolects).
    Level 3: variation mechanisms themselves (transmission rules).
  \item \textbf{Present as $\mathcal{O}_W$-algebra (UNPROVEN).}
    Encode the dynamics of linguistic transmission as wiring-diagram
    operations on linguistic items.
  \item \textbf{Apply $\mathrm{Decomp}$ (UNPROVEN; needs N1).}
    Exhibit the four-slot system.
  \item \textbf{Compute vitality profile (UNPROVEN; needs methodology
    paper).} Verify $k \geq 3$ and $\sigma_i \approx 1$ at every
    level.
\end{enumerate}

\section{Status and follow-up}

\textbf{Current status: UNPROVEN.} The deepest test the framework's
predictive thesis is committed to. Develop after N1 and the
methodology paper land. Expect productive pushback from linguistics
and philosophy of biology.

\bibliographystyle{plain}
\bibliography{references}

\end{document}
